\documentclass[10pt,a4paper]{amsart}
\usepackage[margin=19mm]{geometry}
\usepackage{amsmath,amssymb,mathtools}
\usepackage[scr=rsfs,cal=euler]{mathalfa}
\usepackage{xfrac}
\usepackage[unicode,hidelinks,bookmarks=false]{hyperref}
\newcommand{\ie}{i.e.\ }
\newcommand{\cf}{cf.\ }
\newcommand{\resp}{respectively\ }
\newcommand{\Subsection}{\subsection}
\providecommand{\nobreakdash}{\nobreak\hbox{-}}
\setlength{\emergencystretch}{2em}
\multlinegap0pt
\usepackage{euscript}
\newcounter{sec}
\newcounter{punct}[sec]
\def\punct{\refstepcounter{punct}{\arabic{sec}.\arabic{punct}.  }}
\newtheorem{theorem}{Theorem}[sec]
\newtheorem{proposition}[theorem]{Proposition}
\newtheorem{prop}[theorem]{Proposition}
\newtheorem{lemma}[theorem]{Lemma}
\newtheorem{cor}[theorem]{Corollary}
\newtheorem{corollary}[theorem]{Corollary}
\newtheorem{observation}[theorem]{Observation}
\newtheorem{definition}[theorem]{Definition}
\newtheorem{conjecture}[theorem]{Conjecture}
\newtheorem{question}[theorem]{Question}
\newtheorem{convention}{Convention}[sec]
\def\COUNTERS{\addtocounter{sec}{1}
              \setcounter{punct}{0}
          \setcounter{equation}{0}
          \setcounter{theorem}{0}
          }
\def\SL{\mathrm {SL}}
\def\SU{\mathrm {SU}}
\def\GL{\mathrm  {GL}}
\def\U{\mathrm  U}
\def\OO{\mathrm  O}
\def\Sp{\mathrm  {Sp}}
\def\SO{\mathrm  {SO}}
\def\SOS{\mathrm {SO}^*}
\def\Aff{\mathrm  {Aff}}
\def\PGL{\mathrm  {PGL}}
\def\PU{\mathrm {PU}}
\def\Gr{\mathrm{Gr}}
\def\OSp{\mathrm {OSp}}
\def\AOSp{\mathrm {AOSp}}
\def\fin{\mathrm{fin}}
\def\rest{\mathrm{res}}
\def\GD{\mathbf {GD}}
\def\Mat{\mathrm{Mat}}
\def\Pfaff{\mathrm {Pfaff}}
\def\B{\mathfrak B}
\def\phi{\varphi}
\def\epsilon{\varepsilon}
\def\kappa{\varkappa}
\def\bSp{\mathbf {Sp}}
\def\aff{\mathrm{aff}}
\def\SF{\boldsymbol{\EuScript{S}} \mathbf{F}}
\renewcommand\O{\mathrm{O}}
\def\le{\leqslant}
\def\ge{\geqslant}
\renewcommand{\Re}{\mathop{\rm Re}\nolimits}
\renewcommand{\Im}{\mathop{\rm Im}\nolimits}
\def\la{\langle}
\def\ra{\rangle}
\def\pia{\pi_\downarrow}
\def\super{\mathrm{s}}
\newcommand{\im}{\mathop{\rm im}\nolimits}
\newcommand{\indef}{\mathop{\rm indef}\nolimits}
\newcommand{\dom}{\mathop{\rm dom}\nolimits}
\newcommand{\codim}{\mathop{\rm codim}\nolimits}
\newcommand{\Hom}{\mathop{\rm Hom}\nolimits}
\newcommand{\Ber}{\mathop{\rm Ber}\nolimits}
\newcommand{\graph}{\mathop{\rm graph}\nolimits}
\newcommand{\nul}{\mathop{\mathsf {null}}\nolimits}
\newcommand{\Lagr}{\mathop{\rm Lagr}\nolimits}
\newcommand{\spin}{\mathop{\rm spin}\nolimits}
\newcommand{\Spin}{\mathop{\rm Spin}\nolimits}
\newcommand{\Wei}{\mathop{\rm Wei}\nolimits}
\newcommand{\SpU}{\mathop{\rm SpU}\nolimits}
\newcommand{\OU}{\mathop{\rm OU}\nolimits}
\newcommand{\ASpU}{\mathop{\rm ASpU}\nolimits}
\def\bF{\mathbf F}
\def\cA{\EuScript A}
\def\cB{\EuScript B}
\def\cC{\EuScript C}
\def\cD{\EuScript D}
\def\cE{\EuScript E}
\def\cF{\EuScript F}
\def\cG{\EuScript G}
\def\cH{\EuScript H}
\def\cJ{\EuScript J}
\def\cI{\EuScript I}
\def\cK{\EuScript K}
\def\cL{\EuScript L}
\def\cM{\EuScript M}
\def\cN{\EuScript N}
\def\cO{\EuScript O}
\def\cP{\EuScript P}
\def\cQ{\EuScript Q}
\def\cR{\EuScript R}
\def\cS{\boldsymbol{\EuScript S}}
\def\cT{\EuScript T}
\def\cU{\EuScript U}
\def\cV{\EuScript V}
\def\cW{\EuScript W}
\def\cX{\EuScript X}
\def\cY{\EuScript Y}
\def\cZ{\EuScript Z}
\def\ccF{\boldsymbol {\cF}}
\def\frA{\mathfrak A}
\def\frB{\mathfrak B}
\def\frC{\mathfrak C}
\def\frD{\mathfrak D}
\def\frE{\mathfrak E}
\def\frF{\mathfrak F}
\def\frG{\mathfrak G}
\def\frH{\mathfrak H}
\def\frJ{\mathfrak J}
\def\frK{\mathfrak K}
\def\frL{\mathfrak L}
\def\frM{\mathfrak M}
\def\frN{\mathfrak N}
\def\frO{\mathfrak O}
\def\frP{\mathfrak P}
\def\frQ{\mathfrak Q}
\def\frR{\mathfrak R}
\def\frS{\mathfrak S}
\def\frT{\mathfrak T}
\def\frU{\mathfrak U}
\def\frV{\mathfrak V}
\def\frW{\mathfrak W}
\def\frX{\mathfrak X}
\def\frY{\mathfrak Y}
\def\frZ{\mathfrak Z}
\def\fra{\mathfrak a}
\def\frb{\mathfrak b}
\def\frc{\mathfrak c}
\def\frd{\mathfrak d}
\def\fre{\mathfrak e}
\def\frf{\mathfrak f}
\def\frg{\mathfrak g}
\def\frh{\mathfrak h}
\def\fri{\mathfrak i}
\def\frj{\mathfrak j}
\def\frk{\mathfrak k}
\def\frl{\mathfrak l}
\def\frm{\mathfrak m}
\def\frn{\mathfrak n}
\def\fro{\mathfrak o}
\def\frp{\mathfrak p}
\def\frq{\mathfrak q}
\def\frr{\mathfrak r}
\def\frs{\mathfrak s}
\def\frt{\mathfrak t}
\def\fru{\mathfrak u}
\def\frv{\mathfrak v}
\def\frw{\mathfrak w}
\def\frx{\mathfrak x}
\def\fry{\mathfrak y}
\def\frz{\mathfrak z}
\def\fros{\mathfrak{s}}
\def\bfa{\mathbf a}
\def\bfb{\mathbf b}
\def\bfc{\mathbf c}
\def\bfd{\mathbf d}
\def\bfe{\mathbf e}
\def\bff{\mathbf f}
\def\bfg{\mathbf g}
\def\bfh{\mathbf h}
\def\bfi{\mathbf i}
\def\bfj{\mathbf j}
\def\bfk{\mathbf k}
\def\bfl{\mathbf l}
\def\bfm{\mathbf m}
\def\bfn{\mathbf n}
\def\bfo{\mathbf o}
\def\bfp{\mathbf q}
\def\bfr{\mathbf r}
\def\bfs{\mathbf s}
\def\bft{\mathbf t}
\def\bfu{\mathbf u}
\def\bfv{\mathbf v}
\def\bfw{\mathbf w}
\def\bfx{\mathbf x}
\def\bfy{\mathbf y}
\def\bfz{\mathbf z}
\def\bfA{\mathbf A}
\def\bfB{\mathbf B}
\def\bfC{\mathbf C}
\def\bfD{\mathbf D}
\def\bfE{\mathbf E}
\def\bfF{\mathbf F}
\def\bfG{\mathbf G}
\def\bfH{\mathbf H}
\def\bfI{\mathbf I}
\def\bfJ{\mathbf J}
\def\bfK{\mathbf K}
\def\bfL{\mathbf L}
\def\bfM{\mathbf M}
\def\bfN{\mathbf N}
\def\bfO{\mathbf O}
\def\bfP{\mathbf P}
\def\bfQ{\mathbf Q}
\def\bfR{\mathbf R}
\def\bfS{\mathbf S}
\def\bfT{\mathbf T}
\def\bfU{\mathbf U}
\def\bfV{\mathbf V}
\def\bfW{\mathbf W}
\def\bfX{\mathbf X}
\def\bfY{\mathbf Y}
\def\bfZ{\mathbf Z}
\def\R {{\mathbb R }}
 \def\C {{\mathbb C }}
  \def\Z{{\mathbb Z}}
  \def\H{{\mathbb H}}
\def\K{{\mathbb K}}
\def\N{{\mathbb N}}
\def\Q{{\mathbb Q}}
\def\A{{\mathbb A}}
\def\T{\mathbb T}
\def\bbA{\mathbb A}
\def\bbB{\mathbb B}
\def\bbD{\mathbb D}
\def\bbE{\mathbb E}
\def\bbF{\mathbb F}
\def\bbG{\mathbb G}
\def\bbI{\mathbb I}
\def\bbJ{\mathbb J}
\def\bbL{\mathbb L}
\def\bbM{\mathbb M}
\def\bbN{\mathbb N}
\def\bbO{\mathbb O}
\def\bbP{\mathbb P}
\def\bbQ{\mathbb Q}
\def\bbS{\mathbb S}
\def\bbT{\mathbb T}
\def\bbU{\mathbb U}
\def\bbV{\mathbb V}
\def\bbW{\mathbb W}
\def\bbX{\mathbb X}
\def\bbY{\mathbb Y}
 \def\ov{\overline}
\def\wt{\widetilde}
\def\wh{\widehat}
\renewcommand\emptyset{\varnothing}
\def\arr{\rightrightarrows}
\def\SS{\smallskip}
\def\ev{{\mathrm{even}}}
\def\od{{\mathrm{odd}}}
\def\even{{\mathrm{even}}}
\def\odd{{\mathrm{odd}}}
\def\q{\quad}
\def\F{\mathbf F}
\def\b{\mathfrak b}
\def\RA{\Longrightarrow}
\def\sm{\smallskip}
\def\0{{\ov 0}}
\def\1{{\ov 1}}
\def\mon{{\mathrm{monic}}}
\def\Diff{\mathrm{Diff}}
\def\SDiff{\mathrm{SDiff}}
\def\SCont{\mathrm{SCont}}
\def\vect{\mathfrak{vect}}
\def\ns{\mathfrak{ns}}
\def\sheis{\mathsf{sheis}}
\def\frheis{\mathfrak{sheis}}
\def\free{{\mathfrak{free}}}
\def\vir{\mathfrak{vir}}
\def\cont{\mathfrak{cont}}
\def\osp{\mathfrak{osp}}
\def\DeltA{\boldsymbol{\Delta}}
\begin{document}
\title[Spherical harmonics, operators of multiplication by coordina]{Spherical harmonics, operators of multiplication by coordinates, and infinitesimal conformal transformations}
\author{Yury A. Neretin}
\date{}
\maketitle
\noindent\emph{Source and licence.} Spherical harmonics, operators of multiplication by coordinates, and infinitesimal conformal transformations. Version arXiv:2609.27570v1 (CC BY 4.0). This material is distributed under the Creative Commons Attribution 4.0 International licence, http://creativecommons.org/licenses/by/4.0/, as recorded in the arXiv OAI licence metadata for this exact version and shown on the abstract page https://arxiv.org/abs/2609.27570v1. Full rights notice: licenses/arxiv-2609.27570-CC-BY-4.0.txt.\par\smallskip
\noindent\textit{Editorial scope.} Counted unit: the original \texttt{theorem} environment \texttt{th:2} at source lines 606--638 together with its complete proof at lines 1685--1685; the delivered document keeps source lines 388--1831 so that every referenced label and every citation resolves inside it. Native editable source is the paper's own concatenated TeX; the AMS wrapper is editorial formatting only. No publisher class, upstream script or unrelated artwork is redistributed.
\section{Introduction}

\COUNTERS

{\bf \punct Spherical harmonics.}
We consider the sphere 
$$
x_1^2+x_2^2+x_3^2=1
$$
in $\R^3$ and the action of the  orthogonal
group $\O(3)$ in the space $C^\infty(S^2)$ of
smooth functions on the sphere.  
We get a direct sum of  
 minimal $\O(3)$-invariant subspaces
 (spherical harmonics) 
$$
\bigoplus_{n=0}^\infty \cH_n,\quad  
\text{where $\dim \cH_n={2n+1}$.}
$$

 It is a well-known (and nontrivial) topic of   theory of special functions and classical analysis in $\R^3$. It is also a topic
of representation theory, which at the first glance looks like a toy.

 
Theory of spherical harmonics was exposed in books by Hobson \cite{Hob};
MacRobert \cite{Mac};
Landau, Lifshitz \cite{LL};
Erd\'elyi, Magnus, Oberhettinger, Tricomi \cite{HTF2};
% Chapter 11;
Gelfand, Minlos, Shapiro \cite{GMS}; % 1958
Tychonov,  Samarski \cite{TS}; % Addendum;
Vilenkin \cite{Vil};
Stein, Weiss \cite{SW}; %, Chapter 4; %1971
Perelomov \cite{Per}; % Chapter 4;
Helgason \cite{Hel}; % Introduction, Sect.3;  %1984
Groemer \cite{Gro};
Terras \cite{Ter}; % Chapter 2;
Andrews, Askey, Roy \cite{AAR};
Atkinson, Han \cite{AH}. % Chapter 9.
These expositions overlap but do not cover one another.
We discuss a topic outside these works.

\sm

{\bf\punct The kernel.} 
We define the following function 
\begin{equation}
K(x_1,x_2,x_3; u, n):=
\biggl(\frac{\bigl(1-x_3+u(x_1-ix_2)\bigr)
\bigl(x_1+ix_2+u(x_3-1)\bigr)}{2(1-x_3)}\biggr)^n,
\label{eq:K-x}
\end{equation}
where $(x_1,x_2,x_3)$ ranges in the sphere $S^2$,
$n$ ranges in the space of nonnegative integers, and $u$ ranges in $\C$.

Passing to the spherical coordinates 
$$
x_1=\cos\theta\sin\psi,\quad x_2=\cos\theta \cos\psi,
\quad
x_3=\sin\theta
$$
where $\psi\in [0,2\pi]$, $\theta\in [-\pi/2,\pi/2]$,
we come to the expression
\begin{equation}
K(\psi,\theta;u,n)=
%\Bigl(\frac i2\cdot (u^2+1)\sin 2\theta-2 i u\cos 2\theta+
%i (u^2-1)\sin 2\theta\sin\psi
%\Bigr)^n
\Bigl(\frac 12\cdot \bigl(i e^{-i\psi}\cos\theta
+2u \sin \theta+ i u^2 e^{i\psi}\cos\theta
\bigr)
\Bigr)^n
\label{eq:K-phi}
\end{equation}

Below we prefer to regard $S^2$ as the Riemann sphere
and pass to the stereographic coordinates
$$
z=\frac{x_1+ix_2}{1-x_3}\,\in\C,
$$ 
the kernel in these coordinates has the form
$$
K(z, \ov z; u,n):=\Bigl(\frac{(1+\ov z u)(z-u}{1+z\ov z}\Bigr)^n
.$$

For each $n$ we consider the linear operator
$$
A_n f(u)=\int_{S^2} K(x; u,n) \,f(x)\,d\sigma(x)
$$
from the space of $C^\infty(S^2)$
to the space of polynomials in $u$ of degree
$\le 2n$. By $d\sigma(x)$ we denote the probabilistic
Lebesgue measure on the sphere. 

In the stereographic coordinates these operators have the form
$$
A_n f(u)=\frac 1\pi \int_\C K(z,\ov z;u,n)\, f(z,\ov z)\,\frac{d\Re z\,d\Im z}{(1+z\ov z)^2}.
$$
These operators are intertwining operators from $C^\infty(S^2)$
to $(2n+1)$-dimensional irreducible representations of $\SO(3)$
(see below Proposition \ref{pr:intertwiner}).
We consider all these functions together
regarding the collection $\{\frf_n\}:=\{A_n f\}$, where $n\ge 0$, as a function 
$\frF(u,n)$ depending
on a nonnegative integer $n$ and a complex variable $u$, 
$$
\frF(u,n):=\frf_n(u).
$$
Denote by $\cW^\infty$ the space of functions in $(u,n)$
obtained in this way. Its explicit description is given below
in Subsection \ref{ss:image-C-infty}.



\sm

{\sc Remark.} The intertwining operator $C^\infty(S^2)\to \cH_n$ is defined canonically up to a constant factor depending on $n$, and all ways to write 
this operator leads to the same
result, see for example, \cite{HTF2}, Sect. II.5.1.
\hfill $\boxtimes$


\sm

{\bf \punct The statement of the paper.} We work in the stereographic coordinates $z$, $\ov z$ on the sphere. 
It is easy to see that 
the transformation $f\mapsto \frF$ establishes the following
correspondence of  operators
in $C^\infty(S^2)$ and $\cW^\infty$:
\begin{align}
\frac{\partial}{\partial z}+\ov z^2 \frac{\partial}{\partial \ov z}
\quad&\longleftrightarrow\quad \frac{\partial}{\partial u}
;
\label{eq:triv1}
\\
z  \frac{\partial}{\partial z}- \ov z \frac{\partial}{\partial \ov z}
\quad&\longleftrightarrow\quad
u\frac{\partial}{\partial u}+n
;
\label{eq:triv2}
\\
z^2\frac{\partial}{\partial z}+ \frac{\partial}{\partial \ov z}
\quad&\longleftrightarrow\quad
u^2 \frac{\partial}{\partial u}+2n u
.
\label{eq:triv3}
\end{align} 
This is simply the intertwining property (the operator $A_n$ commutes with the Lie
algebra of the orthogonal group).

Denote by $T_+$ and $T_-$ the shift operators
in $\cW^\infty$:
$$
T_\pm\, \frF(u,n)=\frF(u, n\pm 1),\qquad T_-=T_+^{-1}.
$$ 

\medskip

First, consider the operators of multiplication by the following
functions:
$$
\frac{\ov z}{1+z \ov z},\qquad
\frac{1- z\ov z}{1+z \ov z},\qquad
\frac{z}{1+z \ov z}.
$$
In coordinates $(x_1,x_2,x_3)$ these operators are multiplications
by the functions 
$$
x_1-i x_2,\qquad x_3,\qquad x_1+i x_2.
$$

\begin{theorem}
\label{th:1}
We have the following correspondence of operators
in $C^\infty(S^2)$ and $\cW^\infty$:
\begin{align}
\frac{\ov z}{1+z \ov z}\quad&\longleftrightarrow\quad
-\frac1{2(n+1)(2n+1)}\, T_+\, \frac{\partial^2}{\partial u^2}
+\frac{n}{2(2n+1)}\,T_-
;
\label{eq:drob1}
\\
\frac{1- z\ov z}{1+z \ov z}\quad&\longleftrightarrow\quad
\frac1{(n+1)(2n+1)} \, T_+\, u \frac{\partial^2}{\partial u^2}
+\frac 1{n+1}\, T_+ \frac{\partial}{\partial u}
-\frac{n}{2n+1}\, T_- \,u
;
\label{eq:drob2}
\\
\frac{z}{1+z \ov z} \quad&\longleftrightarrow\quad
-\frac{1}{2(n+1)(2n+1)}\, T_+\, u^2 \frac{\partial^2}{\partial u^2}
+\frac{2}{(n+1)}\, T_+\, u \frac{\partial}{\partial u}-
\label{eq:drob3}
\\
\quad&\phantom{\longleftrightarrow}\quad
\qquad\qquad\qquad\qquad\qquad\qquad\qquad
-2T_+ 
- \frac n{2(2n+1)} \,T_- \, u^2
\notag
.
\end{align}
\end{theorem}

%The conformal groups of $S^2$ is the group $\O(1,3)$ of isometries
%of the 4-dimensional linear space $\R^{1,3}$ equipped with
%the bilinear form
%$$
%\bigl\{ (x_0,x_1,x_2,x_3),(y_0,y_1,y_2,y_3) \bigr\}
%=x_0 y_0+x_1 y_1+x_2 y_2+x_3 y_3
%$$
%It is a 6-dimensional group, its contains 3-dimensional
%group $\O(3)$ of rotations. For elements of the Lie algebra of %rotations
%  the correspondence of operators is given by formulas
%  \eqref{eq:triv1}--\eqref{eq:triv3}.
  
Next, we present operators in $\cW^\infty$ corresponding
to conformal vector fields on the sphere.    


\begin{theorem}
\label{th:2}
We have the following correspondences of operators
in $C^\infty(S^2)$ and $\cW^\infty$:
\begin{equation}
\frac{\partial}{\partial z}- \ov z^2\frac{\partial}{\partial \ov z}
\quad\longleftrightarrow\quad 
-\frac{n + 2}{(n+1)(2 n+1)}\,
T_+ \,\frac{\partial^2}{\partial u^2}  - 
 \frac{n (n -  1)}{(2 n+1)}\, T_-  ;
\label{eq:conf1}
\end{equation}
\begin{multline}
z\frac{\partial}{\partial z}+ \ov z\frac{\partial}{\partial \ov z}
\quad\longleftrightarrow\quad 
\\ - \frac{n+2}{(n+1) (2 n+1)}\,
 T_+\, u \frac{\partial^2}{\partial u^2}+ 
 \frac{n+2}{n+1}\, T_+ \,  \frac{\partial}{\partial u}
 -\frac{(n-1) n}{2 n+1}\, T_- \,u ;
\label{eq:conf2}
\end{multline}
\begin{multline}
z^2\frac{\partial}{\partial z}- \frac{\partial}{\partial \ov z}
\quad\longleftrightarrow\quad 
-\frac{(n+2) }{(n+1) (2 n+1)}\,\, T_+ \, u^2\frac{\partial^2}{\partial u^2}
+\\+
\frac{2 (n+2)}{n+1}\, T_+ \,u  \frac{\partial}{\partial u}
-
2 (n+2)\, T_+
-\frac{n-1}{2 n+1} \,T_-\, u^2.
\label{eq:conf3}
\end{multline}
\end{theorem}


\medskip

Further, denote by $\cS$
the following set of rational functions in $z$, $\ov z$:
\begin{equation}
\Theta^+_{n,k}=\frac{z^k}{(1+z\ov z)^n},\quad \Theta^-_{n,k}:=\frac{\ov z^k}{(1+z\ov z)^n},
\quad\text{where $0\le k\le n$.}
\end{equation}
Denote by $\C[\cS]$ their linear span.

\sm

{\sc Remark.} In coordinates $x_1$, $x_2$, $x_3$, the space $\C[\cS]$
is the space of all polynomials in $x_1$, $x_2$, $x_3$,
see below Observation \ref{obs:3}.
\hfill $\boxtimes$

\begin{corollary}
\label{cor:}
{\rm a)} For any  $0\le k\le n$ and nonnegative integers
$p$, $q$, $r$, $p'$, $q'$, $r'$ 
consider the operator
\begin{equation}
%X^\pm (n,k;p,q,r,p',q',r')=
\Theta_{n,k}^\pm \, \Bigl(\frac{\partial}{\partial z}\Bigr)^p\,
\Bigl(z\frac{\partial}{\partial z}\Bigr)^q \,
\Bigl(z^2\frac{\partial}{\partial z}\Bigr)^r\,
\Bigl(\frac{\partial}{\partial \ov z}\Bigr)^{p'}\,
\Bigl(\ov z\frac{\partial}{\partial \ov z}\Bigr)^{q'} \,
\Bigl(\ov z^2\frac{\partial}{\partial \ov z}\Bigr)^{r'}.
\label{eq:NEBASIS}
\end{equation}
Then the corresponding operator in $\cW^\infty$
is a finite sum
of the form
\begin{equation}
\sum_{M,L\ge 0, K\in \Z} \Xi_{M,L,K}(n) \, T_+^K\, u^M\, \frac{\partial^L}{\partial u^L},
\label{eq:diff-diff}
\end{equation}
where $\Xi_{M,L,K}(n)$ are rational functions in $n$ with poles at
integer and half-integer points.

\sm

{\rm b)} The linear span $\cU$ of operators \eqref{eq:NEBASIS}
is an algebra, i.e., it
is closed with respect to the product.
\end{corollary}

{\sc Remark.}
We have the relation
$$ \Bigl(z\frac{\partial}{\partial z}\Bigr)^2
=
 \Bigl(\frac{\partial}{\partial z}\Bigr) \,
 \Bigl(z^2\frac{\partial}{\partial z}\Bigr)
 -
 z\frac{\partial}{\partial z}
 $$
 and the similar relation for $\frac{\partial}{\partial \ov z}$.
 Therefore, the operators \eqref{eq:NEBASIS} are linearly dependent.
 Moreover, the space $\cU$ is spanned by expressions 
 \eqref{eq:NEBASIS} with $q$, $q'$ being 0 or 1.
 \hfill $\boxtimes$.


\sm


{\bf \punct The general framework.} The statements formulated above are a special case of `operational calculus' for Plancherel decompositions extending operational calculus for the usual Fourier transform on $\R^n$. The phenomenon was observed
in \cite{Ner1} by analysing nonstandard Plancherel formulas for index hypergeometric integral transform \cite{Ner0} and difference equations for
$_3F_2[1]$. A general conjecture was formulated in \cite{Ner1}, \cite{Ner-GL2}.
Several special cases were examined by Molchanov \cite{Mol1}--\cite{Mol3} and the author \cite{Ner1}--\cite{Ner-Lob}. 


Consider the spherical principal series $T_\lambda$ of representations
of the
conformal group $\O(1,3)$ of the sphere (these representations
are realized in a space
of functions on $S^2$, see below Subsection \ref{ss:conformal-sphere}).
We examine restrictions of $T_\lambda$
 to the subgroup $\O(3)$ (this is simply the decomposition
 in spherical harmonics) and write generators
of the Lie algebra
$$
\frs\fro(1,3)_\C=\frs\frl(2,\C)\oplus \frs\frl(2,\C)
$$
 in  decomposition of the restriction. We get the trivial
correspondences \eqref{eq:triv1}--\eqref{eq:triv3}
and also the following correspondences  
 \begin{multline}
\frac{\partial}{\partial z} - \ov z^2 \frac{\partial}{\partial \ov z}
+\frac{2\lambda \ov z}{1+z \ov z} \quad \longleftrightarrow
\\
 -\frac{n +\lambda+ 2}{(n+1)(2 n+1)}\,
T_+ \,\frac{\partial^2}{\partial u^2}  - 
 \frac{n (n -  \lambda - 1)}{2 n+1}\, T_-;
 \label{eq:add1} 
 \end{multline}
 \begin{multline}
z\frac{\partial}{\partial z} + \ov z \frac{\partial}{\partial \ov z}
-\frac{\lambda (1-z \ov z)}{1+z \ov z} \quad \longleftrightarrow
\\ 
- \frac{n+\lambda+2}{(n+1) (2 n+1)}\,
 T_+\, u \frac{\partial^2}{\partial u^2}+ 
 \frac{n+\lambda+2}{n+1}\, T_+ \,  \frac{\partial}{\partial u}
 -\frac{n(n-\lambda-1) }{2 n+1}\, T_- \,u;
\label{eq:add2} 
 \end{multline}
\begin{multline}
z^2\frac{\partial}{\partial z} -  \frac{\partial}{\partial \ov z}
-\frac{2\lambda z}{1+z \ov z} \quad \longleftrightarrow
\\
-\frac{(n+\lambda+2) }{(n+1) (2 n+1)}\, T_+ \, u^2\frac{\partial^2}{\partial u^2}
+\frac{2 (n+\lambda+2)}{n+1}\, T_+ \,u  \frac{\partial}{\partial u}
-
\\-2 (n+\lambda+2)\, T_+
-\frac{n-\lambda-1}{2 n+1} \,T_-\, u^2.
\label{eq:add3}
\end{multline} 

Clearly, these formulas follow from \eqref{eq:drob1}-\eqref{eq:drob3},
\eqref{eq:conf1}-\eqref{eq:conf3} and vise versa. However, it is interesting
that factors depending on $n$ in formulas \eqref{eq:add1}-\eqref{eq:add3}
have a multiplicative structure. A similar picture take place in all  situations
\cite{Mol1}--\cite{Mol3}, \cite{Ner1}--\cite{Ner-Lob} examined earlier.

Our present case is the most simple in the known zoo, but now
 the topic is the most classical. 

\def\skryt 
{\begin{align*}
-\frac{n + 2 \mu}{(1 + n)(1 + 2 n)}\,
T_+ \,\frac{\partial^2}{\partial u^2}  - 
 \frac{n (n - 2 \mu + 1)}{(1 + 2 n)}\, T_-  
 \\
- \frac{2 \mu + n}{(1 + n) (1 + 2 n)}\,
 T_+\, u \frac{\partial^2}{\partial u^2}+ 
 \frac{2 \mu + n}{1 + n}\, T_+ \,  \frac{\partial}{\partial u}
 +\frac{(-1 + 2 \mu - n) n}{1 + 2 n}\, T_- \,u
\end{align*}
\begin{multline*}
-\frac{(2 \mu + n) }{(1 + n) (1 + 2 n)}\, T_+ \, u^2\frac{\partial^2}{\partial u^2}
+\frac{2 (2 \mu + n)}{1 + n}\, T_+ \,u  \frac{\partial}{\partial u}
-
\\-2 (2 \mu + n)\, T_+
+\frac{-1 + 2 \mu - n}{1 + 2 n} \,T_-\, u^2
\end{multline*}}
%\skryt
 
 \sm
 
 {\bf\punct The further structure of the paper.}
 Section 2 contains preliminaries on spherical harmonics.
 Proofs of Theorems \ref{th:1}--\ref{th:2} and Corollary \ref{cor:}
 are contained in Section 3.

\section{Spherical harmonics}

\COUNTERS

{\bf \punct The stereographic projection.} 
We consider the unit sphere $S^2\subset \R^3$
$$
x_1^2+x_2^2+x_3^2=1
$$
equipped with the probabilistic measure $d\sigma$ invariant
with respect to rotations. 
Consider the {\it stereographic projection} of $S^2$ from the north pole
 $(0,0,1)$
to the coordinate plane $x_1 O x_2$. We identify this plane and
$\C$ with the coordinate $z$.
We have
\begin{equation}
z(x)=\frac{x_1+i x_2}{1-x_3}.
\label{eq:z(x)}
\end{equation}
Conversely, 
\begin{equation}
(x_1,x_2,x_3)=\Bigl(\frac{z+\ov z}{1+z\ov z},\,
 \frac{(z-\ov z)}{i(1+z\ov z)},\,
-\frac{1-z\ov z}{1+z\ov z}\Bigr).
\label{eq:x(z)}
\end{equation}

Denote by 
$$
\wh d z:= d\Re z \, d\Im z
$$
the standard Lebesgue measure on $\C$. 
The measure $d\sigma$ on $S^2$ corresponds to the measure
\begin{equation}
\frac1\pi \, \frac{\wh d z}{(1+z \ov z)^{2}}
\label{eq:measure}
\end{equation}
on $\C$.

\sm

{\bf\punct Conformal group of sphere%
\footnote{For a simple self-closed introduction
 to the groups $\SO(3)$, $\SU(2)$,
 $\SL(2,\C)$,
see for example, \cite{KM}, Sections II.11-12, for an introduction to representations of
$\SU(2)$, see for example \cite{FH}, \S\S.10.4, 11.1, see also \cite{GMS}.}.%
\label{ss:conformal-sphere}}
 Consider the Minkowski space $\R^{1,3}$
 with coordinates $(x_0,x_1,x_2,x_3)$ and the indefinite scalar product
$$
%\bigl\{\wt x,\wt y\bigr\}=
\bigl\{ (x_0,x_1,x_2,x_3),(y_0,y_1,y_2,y_3)\bigr\}
=-x_0y_0+x_1y_1+x_2y_2+x_3y_3.
$$
The group of matrices preserving this form is the pseudoorthogonal 
group $\O(1,3)$. It consists of real $(1+3)$-block matrices 
$g=\begin{pmatrix}\alpha&\beta\\\gamma&\delta\end{pmatrix}$ satisfying the condition
$$
g\begin{pmatrix}-1&0\\0&1\end{pmatrix} g^t=\begin{pmatrix}-1&0\\0&1\end{pmatrix}
$$
(here $g^t$ denotes the transposed matrix). The orthogonal group
$\O(3)$ of $\R^3$ is the subgroup in $\O(1,3)$ consisting of matrices
$\begin{pmatrix}1&0\\0&\delta\end{pmatrix}$.
The group $\O(1,3)$, consists of 4  connected components,
 the component $\SO_0(1,3)$ containing the unit is
  determined by two additional conditions:
$$
\alpha>0, \qquad \det \delta>0.
$$

Consider the cone $\cC\subset \R^{1,3}$  
$$
-x_0^2+x_1^2+x_2^2+x_3^2=0
$$
 of isotropic vectors. We identify the section $x_0=1$ of $\cC$
with the sphere $S^2$. Points of this sphere are in one-to-one correspondence
with generatrices of the cone $\cC$.
The group $\O(1,3)$ acts on the set of generatrices and therefore
it acts on $S^2$. Namely, for $x=(x_1,x_2,x_3)\in S^2$
we have
\begin{multline*}
x\mapsto \begin{pmatrix}1&x\end{pmatrix}
\mapsto
 \begin{pmatrix}1&x\end{pmatrix} \begin{pmatrix}\alpha&\beta\\\gamma&\delta\end{pmatrix}
=\begin{pmatrix}\alpha+x\gamma,\beta+x\delta\end{pmatrix}\mapsto
\\
\mapsto 
\begin{pmatrix}1&(\alpha+x\gamma)^{-1}(\beta+x\delta)\end{pmatrix}\mapsto (\alpha+x\gamma)^{-1}(\beta+x\delta)\in S^2.
\end{multline*}
Here we take a point $x$ of $S^2$ and send it to the cone $\cC$;
applying $g$ we get another point   of $\cC$;
intersecting the corresponding generatrix with the hyperplane
$x_0=1$, we get a point of $S^2$.
 
These maps $S^2\to S^2$ are conformal, the coefficient
of dilatation of a transformation $g$ at a point $x$
is 
$$
k(g,x)=
(\alpha+x\gamma)^{-2}.
$$    
Clearly,  the coefficient satisfies the chain rule:
$$
k(g_1 g_2,x)=k(g_1,x)\, k(g_2,xg_1).
$$

{\sc Remark.}
For each $\lambda \in\C$ we  define a representation
of $\SO_0(1,3)$ in the space $C^\infty(S^2)$ by the formula
$$
T_\lambda\begin{pmatrix}
\alpha&\beta\\\gamma&\delta
\end{pmatrix} f(x)=
f\bigl((\alpha+x\gamma)^{-1}(\beta+x\delta)\bigr)\cdot (\alpha+x\gamma)^{-2\lambda}.
$$
By the chain rule,
$$
T_\lambda(g_1)\,T_\lambda(g_2)=T_\lambda(g_1g_2).
$$
The representations $T_\lambda$ are called {\it representations
of spherical principal series}.
If $\Re\lambda=1$, then the representations $T_\lambda$
are unitary in the Hilbert space $L^2(S^2)$.
\hfill $\boxtimes$

\sm

{\bf\punct The conformal group of the Riemann sphere and its Lie algebra.}
We consider the group $\SL(2,\C)$ consisting of complex 
matrices $\begin{pmatrix}a&b\\c&d \end{pmatrix}$
with determinant 1. 
By $\SU(2)$ we denote its subgroup consisting of unitary matrices,
elements of $\SU(2)$ have the form
$$
g=\begin{pmatrix}
a&b\\-\ov b&\ov a
\end{pmatrix}, \qquad\text{where $\det g=|a|^2+|b|^2=1$.}
$$

The group $\SL(2,\C)$ acts on the Riemann sphere 
$$\ov\C:=\C\cup \infty$$
 by conformal transformations
$$
g:\, z\mapsto z^{[g]}:=\frac{b+dz}{a+cz}.
$$
The transformation corresponding to the  matrix 
$\begin{pmatrix}-1&0\\0&-1\end{pmatrix}$
is trivial, so actually we have an action 
of the quotient group $\SL(2,\C)/\Z_2$.

We have
\begin{equation}
\wh d z^{[g]}=\frac{\wh d z}{(a+cz)^2\, \ov{(a+zc)}^2}.
\label{eq:transformation-dz}
\end{equation}



For $h\in \SU(2)$ 
we have
\begin{align}
1+\ov {z^{[h]}} u^{[h]}=\frac{1+\ov z u}{(\ov a- b \ov z)(a-\ov b u)},
\label{eq:1+zu}
\\
z^{[h]}-u^{[h]}=\frac{z-u}{(a-\ov b z)(a-\ov b u)}
\label{eq:z-u}
\end{align}

In particular, this implies that for  $h\in \SU(2)$
we have
\begin{equation}
\frac{\wh d z^{[h]}}{(1+ z^{[h]}\ov{z^{[h]}})^2}=
 \frac{\wh d z}{(1+z\ov z )^2},
\label{eq:whdzh}
\end{equation}
i.e., the transformations $z\mapsto z^{[h]}$ 
with $h\in \SU(2)$ preserve the measure \eqref{eq:measure}.






The stereographic projection $S^2\to\C$ identifies
conformal groups of $\ov\C$ and $S^2$, we have the surjective
homomorphism
$$
\SL(2,\C)\to \SO_0(1,3)
$$
sending
$\SU(2)\to \SO(3)$. The kernel in both cases consists of matrices 
$\begin{pmatrix}\pm 1&0\\0&\pm 1 \end{pmatrix}$.
The group $\SU(2)$ is a two-sheeted covering 
of  $\SO(3)$ and $\SL(2,\C)$ is a two-sheeted covering of  $\SO_0(1,3)$.



\sm 

For instance, we  have the following correspondences of 
one-parametric groups in $\SO(3)$ and in $\SU(2)$:
\begin{align*}
\begin{pmatrix}1&0&0\\
0&\cos 2\phi&\sin2\phi\\
0&-\sin 2\phi&\cos 2\phi
\end{pmatrix}&\longrightarrow
\begin{pmatrix}
\cos\phi& i \sin\phi
\\
i \sin\phi&\cos\phi
\end{pmatrix};
\\
\begin{pmatrix}
\cos2\phi&0&\sin 2\phi\\
0&1&0\\
-\sin 2\phi&0&\cos 2\phi
\end{pmatrix}&
\longrightarrow
\begin{pmatrix}
\cos \phi&\sin\phi\\
-\sin\phi&\cos\phi
\end{pmatrix};
\\
\begin{pmatrix}
\cos 2\phi&\sin 2\phi&0\\
-\sin 2\phi&\cos 2\phi&0\\
0&0&1
\end{pmatrix}
&
\longrightarrow
\begin{pmatrix}
e^{i\phi}&0\\
0&e^{-i\phi}
\end{pmatrix}.
\end{align*}
The corresponding vector fields on $\C$ are
respectively
\begin{gather*}
Q_1=i\Bigl((1-z^2)\frac{\partial}{\partial z}-
(1-\ov z^2)\frac{\partial}{\partial \ov z}\Bigr)
,\qquad
Q_2=(1+z^2)\frac{\partial}{\partial z}+
 (1+\ov z^2)\frac{\partial}{\partial \ov z},\\
Q_3= 2i\Bigl(-z \frac{\partial}{\partial z}+
 \ov z \frac{\partial}{\partial\ov z}\Bigr).
\end{gather*}
These vector fields form a basis in the Lie algebra $\mathfrak{su}(2)$
of the group $\SU(2)$. We prefer to work with the complexification
$\mathfrak{su}(2)_\C\simeq \mathfrak{sl}(2,\C)$ and choose the following basis
in $\mathfrak{su}(2)_\C$:
$$
E_-:=\frac12(Q_2-i Q_1),\qquad E_0:=\frac i2 Q_3,\qquad 
E_+=\frac12(Q_2+i Q_1),
$$
i.e.,
\begin{equation}
E_-:=\frac{\partial}{\partial z}+\ov z^2 \frac{\partial}{\partial \ov z},
\qquad 
E_0:=z\frac{\partial}{\partial z}-\ov z\frac{\partial}{\partial \ov z},\qquad
E_+:=z^2\frac{\partial}{\partial z}+\frac{\partial}{\partial \ov z}.
\label{eq:basis-E}
\end{equation}
They satisfy commutation relations
\begin{equation}
[E_0,E_-]=-E_-,\qquad [E_0,E_+]=E_+,\qquad [E_-,E_+]=2E_0.
\end{equation}

The $\SU(2)$-{\it invariant Laplacian} on $\C$ 
is
\begin{equation}
\DeltA=
\frac12 (E_+E_-+E_-E_+)- E_0^2=
\frac14(Q_1^2+Q_2^2+Q_3^2)=
(1+z\ov z)^2 \,\frac{\partial^2}{\partial z\, \partial \ov z},
\end{equation}
the $\DeltA$ is the usual Laplacian on the sphere $S^2$ written
in the stereographic coordinates.

The complexification of the {\it real} Lie algebra $\mathfrak{sl}(2,\C)$
is
$$
\mathfrak{sl}(2,\C)_\C\simeq \mathfrak{sl}(2,\C)\oplus \mathfrak{sl}(2,\C),
$$
The $\mathfrak{su}(2)_\C\simeq \mathfrak{sl}(2,\C)$ is
the diagonal subalgebra in the direct sum.
We define the following vector fields $\in \mathfrak{sl}(2,\C)_\C$
completing the collection \eqref{eq:basis-E} to the basis
in $\mathfrak{sl}(2,\C)_\C$:
\begin{equation}
F_-:=\frac{\partial}{\partial z}-\ov z^2 \frac{\partial}{\partial \ov z},
\qquad 
F_0:=z\frac{\partial}{\partial z}+\ov z\frac{\partial}{\partial \ov z},\qquad
F_+:=z^2\frac{\partial}{\partial z}-\frac{\partial}{\partial \ov z}.
\label{eq:basis-F}
\end{equation}
They satisfy relations
\begin{gather}
[E_-,F_-]= [E_0,F_0]=[E_+,F_+]=0;
\\
[E_0, F_-]=[E_-,F_0]=-F_-,\qquad [E_0, F_+]=[E_+,F_0]=F_+;
\\
 [E_-,F_+]=[E_+,F_-]=2 F_0
 \\
 [F_0,F_-]=-E_-,\qquad [F_0,F_+]=E_+,\qquad [F_-,F_+]=2E_0
\end{gather}
(so, we have a $\Z_2$-grading).

\sm

{\bf \punct Realizations of irreducible representations of $\SU(2)$.}
Recall a model for irreducible representations of $\SU(2)$
proposed in Berezin \cite{Ber}, \S 5 (see more details in this paper).
Fix $n=0$, 1, 2, \dots. We consider the space $\cB_m$ of
holomorphic functions $p(u)$ on $\C$
satisfying the condition
$$
\int_\C p(u)\,\ov{q(u)}\,(1+u\ov u)^{-m-2}\,\wh d u<\infty
.$$
This  space consists of
  polynomials of degree $\le m$.
We equip $\cB_m$  with the inner product
$$
\bigl\la p,q \bigr\ra_m:=
\frac {m+1}\pi\int_\C p(u)\,\ov{q(u)}\,(1+u\ov u)^{-m-2}\,\wh d u.
$$ 
%The $(m+1)$-dimensional irreducible representation of
The group $\SU(2)$ acts in the $(m+1)$-dimensional
 space $\cB_m$ by the formula
$$
\tau_m\begin{pmatrix}
a&b\\-\ov b&\ov a
\end{pmatrix} p(u)=
p\Bigl(\frac{b+\ov a z}{a-\ov b u}\Bigr)\,(a-\ov b u)^m,
$$
these operators are unitary. The Lie algebra
$\mathfrak{su}(2)_\C$
acts by operators
\begin{equation}
E^{(u)}_-:=\frac{\partial}{\partial u},\qquad
E^{(u)}_0:=u \frac{\partial}{\partial u}+\frac{m}{2},
\qquad E^{(u)}_+:=u^2\frac{\partial}{\partial u}+m u.
\end{equation}



Notice that vectors $1$, $u$, \dots, $u^m\in\cB_m$
are pairwise orthogonal and
\begin{equation}
\|u^k\|_m^2=\frac{k!\, (m-k)!}{m!}.
\label{eq:norms-monomials}
\end{equation}
Indeed,
\begin{multline*}
\|u^k\|^2_m
=\frac{m+1}\pi\int_\C |u|^{2k}(1+|u|^2)^{-m-2}\wh d u
=\frac{2\pi(m+1)}{\pi}\int_0^\infty \frac{r^{2k}\,r\,dr}{(1+r^2)^{m+2}}=
\\=
(m+1)
\int_0^\infty \frac{x^k\,dx}{(1+x)^{m+2}}=(m+1)\,B(k+1,m-k+1)=
\frac{k!\, (m-k)!}{m!}.
\end{multline*}
Here and below we use the following standard beta-integral (see for example \cite{AAR},
formula (1.1.20)):
$$
\int_0^\infty \frac{x^{a-1}dx}{(1+x)^{a+b}}=B(a,b).
$$

{\sc Remark.} Formula \eqref{eq:norms-monomials}
implies that the Hilbert space $\cB_m$ is determined by
the reproducing kernel 
$$L(u,w)=(1+u\ov w)^m$$
on $\C$.
In other words, for any $a\in\C$ consider the function
$\Phi_a(u):=L(u,a)$. Then 
$$\qquad\qquad\qquad\qquad\qquad\qquad\quad
p(a)=\bigl\la  p, \Phi_a \bigr\ra_{\cB_m}
\qquad\qquad\qquad\qquad\qquad\qquad\quad
$$
for any $p\in \cB_m$.
\hfill $\boxtimes$

\sm

{\bf \punct Intertwining operators.}
It is well-known that the representation of 
$\SO(3)$ in $L^2(S^2)$ is a direct sum
$$
L^2(S^1)=\bigoplus_{n=0}^\infty \cH_n
$$
of $(2n+1)$-dimensional irreducible representations 
of $\SO(3)$. Subspaces $\cH_n$ are spaces 
of eigenfunctions $\psi$ of the spherical Laplacian,
$$
\DeltA\psi=-n(n+1)\psi.
$$
The functions
$$
\psi_j^{(n)}:=
E_+^j \frac{\ov z^n}{(1+z \ov z)^n}
=\Bigl(z^2\frac{\partial}{\partial z}+\frac{\partial}{\partial \ov z}\Bigr)^j 
\frac{\ov z^n}{(1+z \ov z)^n},\quad \text{where $j=0$, \dots, $2n$,}
$$
form an orthogonal basis in $\cH_n$.
These functions also are eigenfunctions of $E_0$,
$$
E_0\, \psi_j^{(n)}=(-n+j) \,\psi_j^{(n)}.
$$

 


 The representations of $\SO(3)$ in $\cH_n$ are
equivalent to
the representations $\tau_{2n}$ of $\SU(2)/\Z_2$ defined 
in the previous subsection.

\sm

Denote by $K(z, \ov z; u,n)$
the kernel
\begin{equation}
K(z, \ov z; u,n):=
\Bigl(\frac{(1+\ov z u)(z-u}{1+z\ov z}\Bigr)^n.
\end{equation}

\begin{proposition}
\label{pr:intertwiner}
For a fixed $n$ the operator
\begin{equation}
A_n f(u):= \frac1\pi\int_\C K(z, \ov z; u,n)\, 
f(z,\ov z)\frac{\wh d z}{(1+z\ov z)^2}
\label{eq:A-n}
\end{equation}
is an $\SU(2)$-intertwining operator
$$
L^2\bigl(\C,(1+z\ov z)^{-2} \wh d z\bigr)\to \cB_{2n}
.$$
\end{proposition}

{\sc Proof.} Fix $n$.
First, let us check that the integral is well defined. 
 The expression 
$$
K(\cdot)=(\dots)^n=\sum_{j=0}^{2n}
c_j(z,\ov z)  u^j
$$
is a polynomial in $u$ of degree $2n$
with coefficients of the form
$$
c_j(z,\ov z)= \frac{\kappa_j(z,\ov z)}
{(1+z\ov z)^n}
,$$
where $\kappa_j(z,\ov z)$ are polynomials
of degree $2n$. So, the coefficients $c_j(z,\ov z)$
are bounded functions, they are contained in our $L^2$.
Hence, for $f\in L^2$,
$$
A_n f(u)=\frac1\pi\sum_{j=0}^{2n} u^j\int_\C c_j(z,\ov z)\, f(z,\ov z)
\frac{\wh d z}{(1+z\ov z)^2}=\frac1\pi\sum_{j=0}^{2n} u^j\cdot
\bigl\la\, c_j,\ov {\phantom{.}f \phantom{.}}\, \bigr\ra_{L^2(\dots)}.
$$
The inner products are well-defined,  therefore we get a well-defined operator $L^2\to \cB_{2n}$.



\sm

Keeping in  mind \eqref{eq:1+zu}--\eqref{eq:z-u}, we get
that for $h\in\SU(2)$
\begin{equation}
K\bigl(z^{[h]}, \ov {z^{[h]}};u^{[h]},n\bigr)
=\frac{ K(z,\ov z; u,n)}{(a-\ov b u)^{2n}} 
.
\label{eq:KK}
\end{equation}
Consider the composition of the transformation
$f(z,\ov z)\mapsto f\bigl(z^{[h]}, \ov{z^{[h]}}\bigr)$
and the operator $A_n$.
Keeping in the mind   \eqref{eq:whdzh} and \eqref{eq:KK},
we get
\begin{multline*}
\frac1\pi\int_\C K(z,\ov z;u,n)\, f\bigl(z^{[h]}, \ov{z^{[h]}}\bigr)\,
\frac{\wh d z}{(1+z \ov z)^2}
=\\=
\frac1\pi\int_\C K(z,\ov z;u,n)\,f\bigl(z^{[h]}, \ov{z^{[h]}}\bigr)
\,
\frac{\wh d z^{[h]}}{(1+z^{[h]} \ov {z^{[h]}})^2}
=\\=
\frac1\pi\int_\C 
K\bigl(z^{[h]} ,\ov {z^{[h]}} ;u^{[h]} ,n\bigr)\,
(a-\ov b u)^{2n}\,
f\bigl(z^{[h]}, \ov{z^{[h]}}\bigr)\,
\frac{\wh d z^{[h]}}{(1+z^{[h]} \ov {z^{[h]}})^2}.
\end{multline*}
We substitute $\xi=z^{[h]}$ and come to
$$
\qquad
\biggl(\frac1\pi\int_\C 
K\bigl(\xi ,\ov \xi ;u^{[h]} ,n\bigr)\,
f(\xi, \ov \xi)
\frac{\wh d \xi}{(1+\xi \ov \xi)^2}\biggr)
\cdot (a-\ov b u)^{2n}=
\tau_{2n} \, A_n f(u).
\qquad
\square
$$

\sm 

{\bf \punct The transformation $\boldsymbol{\{A_n\}}$
in the space  $\boldsymbol{L^2(S^2)}$.}

\begin{lemma}
\label{l:norms}
Let $\psi\in \cH_n$.
Then 
$$
\|A_n\psi\|^2_{\cB_{2n}}=\frac{n!\,n!}{(2n+1)!}\, \|\psi\|^2_{\cH_n}.
$$
\end{lemma}

{\sc Proof.}
The operator $A_n:\cH_n \to\cB_{2n}$   intertwines
two irreducible unitary representations; by the Schur lemma,
the ratio $\|A_n\psi\|^2/\|\psi\|^2$ does not depend on $\psi$.
Let us evaluate it for  the function
$\psi_0^{(n)}(z,\ov z)=\ov z^n/(1+z \ov z)^n\in \cH_n$.

We have
\begin{multline*}
\|\psi_0^{(n)}\|^2_{L^2(\C,(1+z\ov z)^{-2}\wh d z)}=
\frac 1\pi \int_\C \frac{\ov z^n}{(1+z\ov z)^n}
\cdot \frac{z^n}{(1+z\ov z)^n}\cdot \frac{\wh d z}{(1+z\ov z)^2} 
=\\=
\frac{2\pi}{\pi}\int_0^{\infty}
\frac{r^{2n}\cdot rdr}{(1+r^2)^{2n+2}}=\int_0^\infty\frac{\rho^n d\rho}{(1+\rho)^{2n+2}}=B(n+1,n+1)=\frac{n!\,n!}{(2n+1)!}.
\end{multline*}
Next,
\begin{multline*}
A_n \psi_0^{(n)}(u)
=\frac 1\pi \int_\C \Bigl(\frac{(1+\ov z u)(z-u}{1+z\ov z}\Bigr)^n
\cdot
  \frac{\ov z^n}{(1+z\ov z)^n}\cdot \frac{\wh d z}{(1+z\ov z)^2} 
  =\\=
 \frac 1\pi\int_0^\infty \int_0^{2\pi}
 \frac{\bigl[1+u r e^{-i\phi} \bigr]^n \, \bigl[r e^{i\phi}-u\bigr]^n\, r^n e^{-i\phi n}\,r\,d\phi\,dr}{(1+r^2)^{2n+2}} 
 =\\=
  \frac 1\pi\int_0^\infty \int_0^{2\pi}
 \frac{\bigl[1+u r e^{-i\phi} \bigr]^n \, \bigl[r -u e^{-i\phi}\bigr]^n\, r^n \,r\,d\phi\,dr}{(1+r^2)^{2n+2}} .
\end{multline*}
We open brackets
$[\dots]^n [\dots]^n$
and get
$$
[\dots]^n [\dots]^n=r^n+\sum_{k=1}^{2n} p_k(u,r) e^{-ik\phi}.
$$
After integration $\int_0^{2\pi}$ all summands except the first one vanish
and we get
$$
2\int_0^\infty\frac{r^n\cdot r^n\,r\,dr}{(1+r^2)^{2n+2}}=
\frac{n!\,n!}{(2n+1)!}= \frac{n!\,n!}{(2n+1)!}\cdot u^0.
$$
By \eqref{eq:norms-monomials}, we have
$$
\|u^0\|^2_{\cB_{2n}}=1.
$$
Therefore,
$$
\|A_m \psi_0^{(n)}\|_{\cB_{2n}}^2=\Bigl(\frac{n!\,n!}{(2n+1)!}\Bigr)^2,
$$
and we come do the desired statement.
\hfill $\square$ 

\sm

Thus, for a function $f$ on the sphere $S^2$
we have a sequence
$$
(\frf_0,\frf_1,\frf_2,\dots)=
(A_0 f,A_1 f, A_2 f,\dots), 
$$
where $\frf_j=A_j f\in \cB_{2n}$.
We regard  such sequences as functions $\frF(u,n)$
depending on a nonnegative integer $n$ and a complex variable
$u$; for a given $k$ the degree 
of a polynomial $\frF(n,u)\Bigl|_{n=k}$ is $\le 2k$.

Clearly, the image of $L^2(S^2)$ 
consists of functions $\frF(u,n)$ satisfying
$$
\sum_{k\ge 0} \frac{(2k+1)!}{k!\,k!}\,\,
\Bigl\| \frF\bigl|_{n=k} \Bigr\|^2_{\cB_{2k}}<\infty.
$$
It
is a Hilbert with inner product defined by 
$$
\bigl\la \frF,\frF'\bigr\ra=\sum_{k=0}^\infty 
 \frac{(2k+1)!}{k!\,k!}\,\,
  \Bigl\la \frF\bigl|_{n=k},\frF'\bigl|_{n=k} \Bigr\ra_{\cB_{2k}}
  .
$$
By the Stirling formula, the asymptotics of coefficients
in this formula is
\begin{equation}
\frac{(2k+1)!}{k!\,k!}
=
(2k+1)\cdot \frac{(2k)!}{k!\,k!}
 \sim 2k\cdot \frac{2^{2k}}{\sqrt {\pi k}}
\qquad
\text {as $k\to+\infty$.}
\label{eq:stirling}
\end{equation}

{\bf \punct The transformation $\boldsymbol{\{A_n\}}$
in the space  $\boldsymbol{C^\infty(S^2)}$.%
\label{ss:image-C-infty}}
A function
$$
F=\sum_{n=0}^\infty \psi_n, \qquad \text{where $\psi_n\in \cH_n$}
$$
is $C^\infty$-smooth if and if for any%
\footnote{The operator $(1-\DeltA)$ acts in
each $\cH_n$ as a multiplication by $1+n(n+1)$.
It is an elliptic positive definite  operator of order 2, therefore its powers
$(1-\DeltA)^{t/2}$, where $t\in \R$, are continuous invertible
pseudo-differential operators from $L^2(S^2)$ to the Sobolev space $W_2^{-t}(S^2)$,
see for example \cite{Tay}, Section II.6.
The space $C^\infty(S^2)$ is $\cap_{t>0} W_2^{t}(S^2)$.} $L>0$
the following estimate holds
$$
\|\psi_n\|_{L^2(S^2)}=o(n^{-L}) \quad \text {as $n\to+\infty$.}
$$
By \eqref{eq:stirling}, the image $\cW^\infty$ of $C^\infty(S^2)$
consits of all functions $\frF(u,n)$ such that for all $L>0$ 
$$
\Bigl\|\frF\bigr|_{n=k} \Bigr\|_{\cB_{2k}}=o(2^{-k}k^{-L})\qquad \text {as $k\to+\infty$.}
$$







%{\bf \punct Spherical representations of 
%$\boldsymbol{\SL(2,\C)}$.}
%For $\lambda\in \C$ we define
%the space $C^\infty_\lambda(\C)$
%consisting of $C^\infty$-functions $f$ on $\C$
%such that
%$$
%f(z^{-1})\cdot |z|^{-2\lambda}
%$$
%also 
%are smooth (this is a condition for asymptotics of $f(z)$ at %infinity).
%We define
% representations $\wt T_\lambda$
% of $\SL_2(\C)$ in $C^\infty_\lambda(\C)$.
%by 
%\begin{equation}
%\wt T_\lambda\begin{pmatrix}a&b\\c&d\end{pmatrix}
%f(z)=f\bigl(z^{[g]}\bigr) |a+zc|^{2\lambda}
%\label{eq:wtT}
%\end{equation}
%(for details, see \cite{GGV}, ???). If $\Re\lambda=1$, then these representations are unitary in $L^2(\C)$. 

%Since we use the measure \eqref{eq:measure},
%we modify  expression \eqref{eq:wtT} in the following way:
%$$
%T'_\lambda\begin{pmatrix}a&b\\c&d\end{pmatrix}
%f(z)=f\bigl(z^{[g]}\bigr) |a+zc|^{2\lambda}\cdot 
%\Bigl(\frac{1+\ov{z^{[g]}} z^{[g]}}{1+\ov z z}\Bigr).
%$$
%The map 
%$$
%f(x)\mapsto f(z(x)) (1+\ov{z(x)} z(x))^{-1}
%$$
%identifies representations $T_\lambda$ given by %\eqref{eq:Tlambda}
%and $T'_\lambda$.

\section{Calculations}

\COUNTERS

{\bf \punct Smoothness of the kernel.}
In the proof of Proposition \ref{pr:intertwiner}
we show that the operators $A_n$ are well-defined. This argument
is not sufficient for differentiation of the integral 
\eqref{eq:A-n} with respect to conformal vector fields.
However, the singularity at $\infty$ is absent, actually
our expression is smooth on the sphere  $S^2$.

To observe the smoothness, we represent the expression in the big bracket
in \eqref{eq:K-x} near the point $x=(0,0,1)$
as
\begin{equation*}
\frac12\Bigl[\frac{u(x_1^2+x_2^2)}{1-x_3}+ \Bigl(-u^2(x_1-i x_2)+(x_1+i x_2)\Bigr)
-(1-x_3) u\Bigr].
\end{equation*}
The second and the third summands in the brackets $[\dots]$
 are smooth on $S^2$. 
In the first summand, we
eliminate the irrationality in the denominator:
$$
\frac{u(x_1^2+x_2^2)(1+x_3)}{(1-x_3)(1+x_3)}=
\frac{u(x_1^2+x_2^2)(1+x_3)}{x_1^2+x_2^2}= u(1+x_3)
,
$$
and we again get a smooth expression.

\sm

{\bf \punct The verification of Theorem \ref{th:1}.}
First, we verify \eqref{eq:drob1}. Denote by 
$$
\cL:=-\frac1{2(n+1)(2n+1)}\, T_+\, \frac{\partial^2}{\partial u^2}
+\frac{n}{2(2n+1)}\,T_-
$$
 the operator
in the right side of \eqref{eq:drob1}. Our statement
is equivalent to the identity
$$
\cL \,K(z,\ov z; u, n)=\frac{\ov z}{1+z\ov z}\cdot K(z,\ov z; u, n).
$$
Now it can be verified by any system of computer algebra%
\footnote{The author used Wolfram Mathematica.}. We wish to explain how to find this identity and how to verify it by hands.

Let us transform to a convenient for us form the expressions
$$
T_- K,\qquad \frac{\partial}{\partial u} T_+,
\qquad \frac{\partial^2}{\partial u^2} T_+.
$$

Denote
$$
\kappa:=\frac{(1+\ov z u)(z-u)}{1+z\ov z},
$$
so $K=\kappa^n$.
Notice that 
$$
\kappa^{-1}=
\frac{1+z\ov z}{(1+\ov z u)(z-u)}=\frac{\ov z}{1+ \ov z u}+
\frac1{z-u}.
$$
Therefore,
\begin{equation}
T_- K=\Bigl(\frac{\ov z}{1+ \ov z u}+
\frac1{z-u}\Bigr) \cdot K.
\label{eq:T-K}
\end{equation}

Applying the formula
\begin{equation}
\frac{\partial}{\partial x}\prod_j f_j(x)^{m_j} 
=
\Bigl(\sum_j \frac{m\,  f'_j(x) }{f_j(x)} \Bigr)
\cdot
\prod_j f_j(x)^{m_j}, 
\label{eq:logarithmic}
\end{equation}
we get
\begin{multline}
\frac{\partial^2}{\partial u^2}\, T_+ K
=\frac{\partial^2}{\partial u^2} \kappa^{n+1}=
\frac{\partial}{\partial u}\Bigl[(n+1)\Bigl(\frac{\ov z}{1+\ov z u}-\frac{1}{z-u} \Bigr) \kappa^{n+1}\Bigr]
=\\=
\Bigl\{(n+1)^2\Bigl(\frac{\ov z}{1+\ov z u}-\frac{1}{z-u} \Bigr)^2
-(n+1)\Bigl(\frac{\ov z^2}{(1+\ov z u)^2}+\frac{1}{(z-u)^2} \Bigr)\Bigl\}
\,\kappa^{n+1}
=\\=
\Bigl\{n(n+1)\Bigl(\frac{\ov z^2}{(1+\ov z u)^2}+\frac{1}{(z-u)^2} \Bigr)-2(n+1)^2 \frac{\ov z}{(1+\ov z u)(z-u)}
\Bigr\}\kappa\cdot \kappa^n.
\end{multline}
We transform summands in the right hand side,
$$
\frac{\ov z}{(1+\ov z u)(z-u)}\cdot \kappa=\frac{\ov z}{1+z \ov z};
$$
$$
\frac{\ov z^2}{(1+\ov z u)^2}\cdot \kappa=
\frac{\ov z^2(z-u)}{(1+\ov z u)(1+\ov z z)}
=\frac{\ov z}{1+\ov z u} - \frac{\ov z}{1+\ov z z};
$$
$$
\frac{1}{(z-u)^2}\cdot \kappa =\frac{1+\ov z u}{(z-u)(1+\ov z z)}
=\frac{1}{z-u}-\frac{\ov z}{1+\ov z z}.
$$
Thus, we have
\begin{equation}
\frac{\partial^2}{\partial u^2}\, T_+ K
=\Bigl\{
n(n+1)\Bigl(\frac{\ov z}{1+\ov z u}+ \frac{1}{z-u}\Bigr)
+ \frac{C\ov z}{1+z\ov z}
\Bigr\}\cdot K,
\label{eq:partial2K}
\end{equation}
where $C=-2n(n+1)-2(n+1)^2$. Comparing the last expression
with \eqref{eq:T-K} we observe that
$$
u^2\frac{\partial^2}{\partial u^2}\,T_+\, K-2n(n+1) \,T_- K=
-2(n+1)(2n+1)\cdot \frac{\ov z}{1+z \ov z}\cdot K.
$$
This establishes \eqref{eq:drob1}.

\sm

Next,
$$
\Bigl[z^2\frac{\partial}{\partial z}+\frac{\partial}{\partial \ov z},\,\,
\frac{\ov z}{1+z \ov z}\Bigr]=\frac{1 - z\ov z}{1+z \ov z}.
$$
We know the operator corresponding $z^2\frac{\partial}{\partial z}+\frac{\partial}{\partial \ov z}$  (see \eqref{eq:triv3}),
and we obtained the operator $\cL$ corresponding to $\frac{\ov z}{1+z \ov z}$
(see \eqref{eq:triv3}). Evaluating their commutator
$$
\Bigl[u^2 \frac{\partial}{\partial u}+2nu,\,\,
-\frac1{2(n+1)(2n+1)}\, T_+\, \frac{\partial^2}{\partial u^2}
+\frac{n}{2(2n+1)}\,T_-\Bigr],
$$
we come to  formula \eqref{eq:drob2}.

\sm

Repeating the same reasoning with the commutator
$$
\Bigl[z^2\frac{\partial}{\partial z}+\frac{\partial}{\partial \ov z},
\,\,
\frac{1-z\ov z}{1+z \ov z}\Bigr]=-\frac{2 z}{1+z \ov z},
$$
after straightforward calculations
come to \eqref{eq:drob3}.

\sm

{\bf \punct Integration by parts.}
We say that an operator $Q$ in the space of functions
on $\C$ is {\it formally transposed} to an operator $R$
if
$$
\int_\C  \bigl(Qf(z,\ov z)\bigr)\cdot g(z,\ov z)
\,(1+z \ov z)^{-2}  \wh d z=
\int_\C  f(z,\ov z)\cdot \bigl(Rg(z,\ov z)\bigr) \,
(1+z \ov z)^{-2}  \wh d z
$$
for all $f$, $g$, which are actually are smooth on the sphere $S^2$.
Integrating by parts, 
\begin{multline}
\int_\C  \Bigl(\frac{\partial}{\partial z} f(z)\Bigr) g(z,\ov z)
\frac{\wh d z}{(1+z \ov z)^2}
=-\int_\C f(z,\ov z) \cdot \frac{\partial}{\partial z}\bigl( g(z,\ov z)
(1+z \ov z)^{-2} \bigr) \wh d z
=\\=
-\int_\C f(z,\ov z) \cdot \Bigl(\frac{\partial}{\partial z}-
\frac{2\ov z}{1+z\ov z}\Bigr)\, g(z,\ov z)\cdot
\frac{\wh d z}{(1+z \ov z)^{2}}
,
\end{multline}
we get
$$
\Bigl(\frac{\partial}{\partial z}\Bigr)^t
=-\frac{\partial}{\partial z}+ \frac{2\ov z}{1+z \ov z},
\qquad \Bigl(\frac{\partial}{\partial \ov z}\Bigr)^t
=-\frac{\partial}{\partial \ov z}+ \frac{2 z}{1+z \ov z}.
$$
For operators \eqref{eq:basis-E},
we have 
$$
E_\pm^t=-E_\pm^t, \qquad E_0^t=-E_0.
$$
For operators
\eqref{eq:basis-F}, straightforward calculations show
\begin{equation}
F_-^t=\Bigl(\frac{\partial}{\partial z}-
\ov z^2 \frac{\partial}{\partial \ov z} \Bigr)^t
=-F_-+\frac{4 \ov z}{1+z\ov z};
\label{eq:FF}
\end{equation}
$$
F_0^t=\Bigl(z\frac{\partial}{\partial z}+
\ov z \frac{\partial}{\partial \ov z} \Bigr)^t=
-F_0-2\frac{1-z \ov z}{1+z\ov z};
$$
$$
F_+^t=\Bigl(z^2\frac{\partial}{\partial z}-
 \frac{\partial}{\partial \ov z} \Bigr)^t=
-F_+-\frac{4z}{1+z\ov z}.
$$


{\bf \punct The proof of Theorem \ref{th:2}.}

Let us verify the statement \eqref{eq:conf1}.
By \eqref{eq:FF}, we get
\begin{multline*}
\int_\C K(z,\ov z;u,n)\, F_- g(z,\ov z)\, \frac{\wh d z}{(1+z\ov z)^2}
=\\=
\int_\C\Bigl[\Bigl(-F_-+\frac{4\ov z}{1+z\ov z}\Bigr)
K(z,\ov z;u,n)\Bigr]\,  g(z,\ov z)\, \frac{\wh d z}{(1+z\ov z)^2}.
\end{multline*}
So, we need a differential-difference  operator $\cR$
such that
$$
\Bigl(-F_-+\frac{4\ov z}{1+z\ov z}\Bigr)
K(z,\ov z;u,n) =\cR\, K(z,\ov z;u,n).
$$
By \eqref{eq:drob1}, it is sufficient to find $\cR'$
such that
$$-F_-\,K(z,\ov z;u,n) =\cR'\, K(z,\ov z;u,n).$$
It  to verify that
\begin{equation}
F_-\, K=n\Bigl(\frac{\ov z}{1+\ov z u}+\frac1{z-u}-
\frac{\ov z}{1+z\ov z}\Bigr)\, K.
\label{eq:FK}
\end{equation}
This identity can be found by manipulations with the formula 
\eqref{eq:logarithmic}.

We compare expressions
for
\begin{equation}
K^{-1}\cdot T_- K,\qquad K^{-1}\cdot \frac{\partial^2}{\partial u^2}T_+K,
\qquad K^{-1}\cdot F_-\, K
,
\label{eq:expressions}
\end{equation}
see \eqref{eq:T-K}, \eqref{eq:partial2K}, \eqref{eq:FK}.
Each expression has form
$$
A_j \Bigl(\frac{\ov z}{1+\ov z u}+\frac1{z-u}\Bigr)+ B_j\cdot \frac{\ov z}{1+z\ov z},
$$
where $A_j$, $B_j$ are polynomials in $n$. Therefore 
\eqref{eq:expressions} are linearly dependent with coefficients depending on $n$.
This allows to express $F_-\, K$ as a linear combination
of $T_- K$ and $\frac{\partial^2}{\partial u^2}T_+K$.


  % Now after straightforward calculations we get \eqref{eq:conf1}.

Next, $[E_+,F_-]=2 F_0$. Since we know corresponding
$E_+$ and $F_-$, after straightforward calculations we  evaluate the operator \eqref{eq:conf2} corresponding to $F_0$.

Finally, $[E_+,F_0]=F_+$, and this allow to evaluate \eqref{eq:conf3}.
 
\sm

{\bf \punct Proof of Corollary \ref{cor:}.%
\label{ss:proof-corollary}}
We begin with three simple observations.

\begin{observation}
\label{obs:1}
The space  $\C[\cS]$
is closed with respect to the product.
\end{observation}


%{\sc Proof.}
Clearly, $\Theta^+_{m,k}\, \Theta^+_{n,l}=\Theta^+_{m+n,k+l}$.
Let us examine products
$\Theta^+_{m,k}\cdot \Theta^-_{n,l}$.
 To be definite, assume
that $k\ge l$. Then
\begin{multline*}
\Theta^+_{k,m}\cdot \Theta^-_{l,n}=
\frac{(z\ov z)^l z^{k-l}}{(1+z \ov z)^{m+n}}=
\frac{\bigl[(1+z\ov z)-1\bigr]^l z^{k-l}}{(1+z \ov z)^{m+n}}
=\\=
 \frac{\sum_{j=0}^l C_l^j(-1)^{j}(1+z\ov z)^j\cdot z^{k-l}}
{(1+z \ov z)^{m+n}}=
\sum_{j=0}^l C_l^j(-1)^{j} \Theta^+_{m+n-j,k-l}.
\end{multline*}



\begin{observation}
\label{obs:2}
 The space $\C[\cS]$ is invariant with respect to  the operators
$\frac{\partial}{\partial z}$, $z\frac{\partial}{\partial z}$, 
$z^2\frac{\partial}{\partial z}$, $\frac{\partial}{\partial \ov z}$, $\ov z\frac{\partial}{\partial \ov z}$, 
$\ov z^2\frac{\partial}{\partial \ov z}$.
\end{observation}

\sm

This follows from the following trivial calculations 
(but their result is not a priory obvious). Let $0\le m \le n$. We apply our operators
to $\Theta^+_{n,n-m}$:
$$
z^2 \frac{\partial}{\partial z} \frac{z^{n-m}} {(1+z \ov z)^n}
= \frac{n z^{n-m+1}}{(1+z\ov z)^{n+1}}-
\frac{m z^{n-m+1}}{(1+z\ov z)^{n}}\,\,
=\,\, n\Theta^+_{n+1,n-m+1}-m\Theta^+_{n,n-m+1}.
$$
If $m=0$, we can not write $\Theta^+_{n,n-m+1}$,
but in this case the term with $\Theta^+_{n,n-m+1}$
is absent.
 Next,
$$
z \frac{\partial}{\partial z} \frac{z^{n-m}} {(1+z \ov z)^n}
\!=\! \frac{n z^{n-m}}{(1+z\ov z)^{n+1}}-
\frac{m z^{n-m}}{(1+z\ov z)^{n}},
\,\, 
 \frac{\partial}{\partial z} \frac{z^{n-m}} {(1+z \ov z)^n}
\!=\! \frac{n z^{n-m-1}}{(1+z\ov z)^{n+1}}-
\frac{m z^{n-m-1}}{(1+z\ov z)^{n}}.
$$
The right hand side of the second equality does not have the desired form if $m=n$. But in that case $\frac{\partial}{\partial z} (1+z \ov z)^{-n}=-n \ov z (1+z \ov z)^{-n-1}$. Further, 
$$
\frac{\partial}{\partial \ov z}\,
\frac{z^{n-m}}
{(1+z \ov z)^{n}}=\frac{n\,z^{n-m+1}}
{(1+z \ov z)^{n+1}},\quad 
\ov z\frac{\partial}{\partial \ov z}\,
\frac{z^{n-m}}
{(1+z \ov z)^{n}}=\frac{n\,z^{n-m}}{(1+z \ov z)^{n+1}}-\frac{n\,z^{n-m}}{(1+z \ov z)^{n}};
$$
$$
\ov z^2\frac{\partial}{\partial \ov z}\,
\frac{z^{n-m}}
{(1+z \ov z)^{n}}=-\frac{n z^{n-m-1}}{(1+z \ov z)^{n+1}}+
\frac{2n z^{n-m-1}}{(1+z \ov z)^{n}}-
\frac{n z^{n-m-1}}{(1+z \ov z)^{n-1}}.
$$
The last transformation is not appropriate for us
if $m=n$. In this case,
$$
\ov z^2\frac{\partial}{\partial \ov z}\,
\frac{1}
{(1+z \ov z)^{n}}=-\frac{n z \ov z^2}{(1+z \ov z)^{n+1}}
=\frac{n \ov z}{(1+z \ov z)^{n+1}}- \frac{n \ov z}{(1+z \ov z)^{n}}.
$$ 

\begin{observation}
\label{obs:3}
 In the coordinates $x_1$, $x_2$, $x_3$ the space $\C[\cS]$ is the space of all polynomials. 
\end{observation}


\begin{thebibliography}{99}
\bibitem[AAR]{AAR} Andrews G. E., Askey, R., Roy, R. {\it Special functions.} Cambridge University Press, Cambridge, 1999.
\bibitem[AH]{AH} Atkinson K., Han W. {\it Spherical harmonics and approximations on the unit sphere: an introduction.} Lecture Notes in Math., 2044, Springer, Heidelberg, 2012.
\bibitem[Ber]{Ber} Berezin, F. A. {\it General concept of quantization.} Comm. Math. Phys. 40 (1975), 153-174.
\bibitem[FH]{FH} Fulton W., Harris J. {\it Representation theory. A first course.} Springer-Verlag, New York, 1991.
\bibitem[GMS]{GMS} Gelfand I. M., Minlos R. A., Shapiro Z. Ja., {\it Representations of the rotation group and of the Lorentz group, and their applications}. Pergamon Press, New York, 1963. % (Russian), %Gosudarstv. Izdat. Fiz.-Mat. Lit., Moscow, 1958.
\bibitem[Gro]{Gro} Groemer H. {\it Geometric applications of Fourier series and spherical harmonics.} Cambridge University Press, Cambridge, 1996.
\bibitem[HTF2]{HTF2} Erd\'elyi A., Magnus W., Oberhettinger F., Tricomi F. G. {\it Higher transcendental functions. Vol. II.} Based, in part, on notes left by Harry Bateman. McGraw-Hill Book Co., New York--Toronto--London, 1953.
\bibitem[Hel]{Hel} Helgason S. {\it Groups and geometric analysis. Integral geometry, invariant differential operators, and spherical functions.} Academic Press, Orlando, 1984.
\bibitem[Hob]{Hob} Hobson E. W. {\it The theory of spherical and ellipsoidal harmonics.} Cambridge Univ. Press, Cambridge, 1931.
\bibitem[KM]{KM} Kostrikin A. I., Manin Yu. I. {\it Linear algebra and geometry.} Gordon and Breach, New York, 1989
\bibitem[LL]{LL} Landau L. D., Lifshitz E. M. {\it Quantum mechanics: non-relativistic theory.} Pergamon Press, London--Paris, 1958.
\bibitem[Mac]{Mac} MacRobert Th. M. {\it Spherical harmonics: an elementary treatise on harmonic functions with applications.} Methuen \& Co., London, 1948.
\bibitem[Mol1]{Mol1} Molchanov V. F., {\it Canonical representations and overgroups for hyperboloids.} Funct. Anal. Appl., 39:4 (2005), 284-295.
\bibitem[Mol3]{Mol3} Molchanov V. F., {\it Poisson and Fourier transforms for tensor products}. Funct. Anal. Appl., 49:4 (2015), 279-288.
\bibitem[Ner-GL2]{Ner-GL2} Neretin Yu. A., {\it Operational calculus for the Fourier transform on the group $\GL(2,\R)$ and the problem about the action of an overalgebra in the Plancherel decomposition}, Funct. Anal. Appl., 52:3 (2018), 194-202.
\bibitem[Ner-Lob]{Ner-Lob} Neretin Yu. A. {\it The Fourier transform on the Lobachevsky plane and operational calculus.} Funct. Anal. Appl. 54 (2020), no. 4, 278-286.
\bibitem[Ner0]{Ner0} Neretin Yu. A. {\it The index hypergeometric transform and an imitation of the analysis of Berezin kernels on hyperbolic spaces.} Sb. Math. 192 (2001), no. 3-4, 403-432.
\bibitem[Ner1]{Ner1} Neretin Yu. A., {\it The action of an overalgebra on the Plancherel decomposition and shift operators in the imaginary direction}, Izv. Math., 66:5 (2002), 1035-1046. %Yu. A. Neretin, ``Difference Sturm�Liouville problems in the imaginary direction", J. Spectr. Theory, 3:3 (2013), 237�269. MR3073412
\bibitem[Per]{Per} Perelomov A. {\it Generalized coherent states and their applications.} Springer-Verlag, Berlin, 1986.
\bibitem[SW]{SW} Stein E. M., Weiss G. {\it Introduction to Fourier analysis on Euclidean spaces.} Princeton University Press, Princeton, 1971.
\bibitem[TS]{TS} Tychonov A. N., Samarski A. A. {\it Partial differential equations of mathematical physics.} Vol. II. Holden-Day, San Francisco, 1967.
\bibitem[Tay]{Tay} Taylor M. E. {\it Pseudodifferential operators.} Princeton University Press, Princeton, 1981.
\bibitem[Ter]{Ter} Terras A. {\it Harmonic analysis on symmetric spaces --- Euclidean space, the sphere, and the Poincar\'e upper half-plane.} Second edition. Springer, New York, 2013.
\bibitem[Vil]{Vil} Vilenkin N. Ja. {\it Special functions and the theory of group representations.} American Mathematical Society, Providence, 1968.
\end{thebibliography}
\end{document}
